% Integration fragment; all labels use the urc: prefix.
% Requires amsmath, amssymb, amsthm, hyperref. No custom global macros.
% Bibliography: cite the existing DLMF and Lattice Contacts entries.

\section{Unequal-rate Stieltjes boundary collisions}
\label{urc:section}

The independent-shift product generator in the incoming \emph{Lattice
Contacts and Mixed Zeta Identities} treats positive shifts and proves
compatibility with the coincident finite part along the common-shift
diagonal. Its Question~Q9, repeated at the end of its independent-shift
section \cite[Section~6 and Question~Q9]{LatticeIncoming}, leaves open the path
\[
 a=\lambda\epsilon,\qquad b=\mu\epsilon,\qquad
 \lambda,\mu>0,\qquad \epsilon\downarrow0.
\]
We determine every divergent power and logarithm on this path and its
remaining constant, for arbitrary Stieltjes indices and argument
derivatives. The constant correction is considerably smaller than the
full independent-shift product generator: it depends only on two Taylor
coefficients at one and a finite polynomial in the shift ratio.

These statements resolve the specified collision question relative to
the inspected repository sources. The underlying local-subtraction
method is classical; no claim of historical priority for general
finite-part asymptotics is intended.

\subsection{The coordinate and the finite contact polynomial}

We use the convention
\begin{equation}\label{urc:stieltjes}
 Z(u,z)=\zeta(1+u,z)-\frac1u
       =\sum_{m\ge0}\frac{(-1)^m}{m!}\gamma_m(z)u^m.
\end{equation}
In particular, $\gamma_0(z)=-\psi(z)$. Write
\begin{equation}\label{urc:integral}
 I_{m,n}^{r,t}(\epsilon;\lambda,\mu)
 =\int_0^1\gamma_m^{(r)}(x+\lambda\epsilon)
             \gamma_n^{(t)}(x+\mu\epsilon)\,dx,
 \qquad m,n,r,t\in\mathbb Z_{\ge0}.
\end{equation}
Every integral in \eqref{urc:integral} is ordinary. Its asymptotic
constant means the coefficient of $\epsilon^0(\log\epsilon)^0$.
The inherited coincident finite part is
\begin{equation}\label{urc:finitepart}
 F_{m,n}^{r,t}
 =\operatorname{CT}_{\delta\downarrow0}
    \int_\delta^1\gamma_m^{(r)}(x)\gamma_n^{(t)}(x)\,dx.
\end{equation}
Thus the two constants use the displayed coordinates $\epsilon$ and
$\delta$; no subtraction scale is left implicit.

For $r\ge0$ define a holomorphic germ at $u=0$ by
\begin{equation}\label{urc:contact-germ}
 C_r(u;\lambda,\mu)=\frac{(-1)^r(1+u)_r}{r!}
 \left\{\frac{\lambda^{-u}-1}{u}
 +\lambda^{-u}\sum_{j=1}^r
    \binom rj\frac{(\mu/\lambda-1)^j}{u+j}\right\},
\end{equation}
and put
\begin{equation}\label{urc:kappa}
 \kappa_{m,r}(\lambda,\mu)
 =(-1)^m\left.\partial_u^m C_r(u;\lambda,\mu)\right|_{u=0}.
\end{equation}
The sum is empty if $r=0$. The quotient at $u=0$ is removable.
Moreover, $(1+u)_r/(u+j)$ is a polynomial for $1\le j\le r$,
so \eqref{urc:contact-germ} contains no hidden special-function limit.

\begin{theorem}[Complete constant correction for unequal rates]
\label{urc:constant-theorem}
For all indices in \eqref{urc:integral},
\begin{equation}\label{urc:constant-formula}
 \boxed{\begin{aligned}
 \operatorname{CT}_{\epsilon\downarrow0}
       I_{m,n}^{r,t}(\epsilon;\lambda,\mu)
 ={}&F_{m,n}^{r,t}
    +\gamma_n^{(r+t)}(1)\kappa_{m,r}(\lambda,\mu)\\
   &+\gamma_m^{(r+t)}(1)\kappa_{n,t}(\mu,\lambda).
 \end{aligned}}
\end{equation}
In particular, with $R=\mu/\lambda$,
\begin{equation}\label{urc:zeroth-contact}
 \boxed{\displaystyle
 \kappa_{0,r}(\lambda,\mu)
 =(-1)^r\left\{-\log\lambda+
                   \sum_{j=1}^r\frac{R^j-1}{j}\right\}.}
\end{equation}
For arbitrary Stieltjes indices without argument differentiation,
\begin{equation}\label{urc:no-argument-derivatives}
 \operatorname{CT}_{\epsilon\downarrow0}I_{m,n}^{0,0}
 =F_{m,n}^{0,0}
  -\frac{\gamma_n(1)}{m+1}(\log\lambda)^{m+1}
  -\frac{\gamma_m(1)}{n+1}(\log\mu)^{n+1}.
\end{equation}
\end{theorem}

Only the Taylor coefficient of degree $r$ in the analytic part of
the second factor and the coefficient of degree $t$ in the analytic
part of the first factor enter \eqref{urc:constant-formula}. In
particular, the singular--singular product changes the divergent
terms but contributes no path-dependent constant.

\subsection{A finite formula for every contact coefficient}

For a fully finite harmonic description, write
\[
 A_r(u)=\frac{(1+u)_r}{r!}
       =\sum_{a=0}^r e_{r,a}u^a,
 \qquad
 Q_r(u;R)=\sum_{j=1}^r\binom rj(R-1)^j\frac{A_r(u)}{u+j}
       =\sum_{a=0}^{r-1}q_{r,a}(R)u^a.
\]
The $Q_0$ sum is empty. The coefficients $e_{r,a}$ are elementary
symmetric polynomials in $1,1/2,\ldots,1/r$, equivalently
\begin{equation}\label{urc:harmonic-generator}
 A_r(u)=\exp\left\{\sum_{h\ge1}
        \frac{(-1)^{h+1}}h H_r^{(h)}u^h\right\}.
\end{equation}
With $L=\log\lambda$, coefficient extraction in
\eqref{urc:contact-germ} gives
\begin{equation}\label{urc:finite-coefficients}
 \begin{aligned}
 \kappa_{m,r}={}&(-1)^{m+r}m!
 \sum_{a=0}^{\min(r,m)}e_{r,a}
          \frac{(-L)^{m-a+1}}{(m-a+1)!}\\
 &+(-1)^{m+r}m!
 \sum_{a=0}^{\min(r-1,m)}q_{r,a}(R)
          \frac{(-L)^{m-a}}{(m-a)!}.
 \end{aligned}
\end{equation}
The second sum is interpreted as empty at $r=0$. For example,
\begin{align*}
 \kappa_{0,0}&=-L,&
 \kappa_{1,0}&=-\frac{L^2}{2},\\
 \kappa_{0,1}&=L+1-R,&
 \kappa_{1,1}&=\frac{L^2}{2}-RL,\\
 \kappa_{0,2}&=-L+R+\frac{R^2}{2}-\frac32.
\end{align*}
At $\lambda=1$ the contact is a polynomial of degree at most $r-1$
in $u$. Consequently $\kappa_{m,r}(1,\mu)=0$ whenever $m\ge r$.
At $\lambda=\mu=1$ all contacts vanish. These are exact truncations
of the correction, not asymptotic statements about its size.

\subsection{The complete singular power--logarithm expansion}

The following evaluated germ specifies every divergent term as well
as the constant. Put
\begin{equation}\label{urc:local-decomposition}
 P_r(u)=(-1)^r(1+u)_r,\qquad
 Z_r(u,x)=\partial_x^r Z(u,x)
         =P_r(u)x^{-r-1-u}+A_r^{\mathrm{reg}}(u,x),
\end{equation}
where $A_r^{\mathrm{reg}}(u,x)=\partial_x^r Z(u,1+x)$.
The superscript distinguishes this analytic function from the finite
polynomial $A_r$ above. Define
\begin{equation}\label{urc:analytic-coefficients}
 a_{t,k}(v)=\frac1{k!}\partial_z^{t+k}Z(v,z)\big|_{z=1}.
\end{equation}
These are holomorphic near $v=0$, and
\[
 (-1)^n\partial_v^n a_{t,k}(0)
      =\frac{\gamma_n^{(t+k)}(1)}{k!}.
\]

The only nonpolynomial local kernel needed is the elementary beta
integral
\begin{align}
 B_{r,t}(u,v;\lambda,\mu)
 &=P_r(u)P_t(v)\int_0^\infty
 \frac{dy}{(y+\lambda)^{r+1+u}(y+\mu)^{t+1+v}}\notag\\
 &=\frac{P_r(u)P_t(v)\mu^{-r-t-1-u-v}}{r+t+1+u+v}\notag\\
 &\quad\cdot{}_2F_1\left(
 \begin{matrix}r+1+u,\ r+t+1+u+v\\r+t+2+u+v\end{matrix};
                   1-\frac\lambda\mu\right).
 \label{urc:beta-kernel}
\end{align}
The integral fixes its branch. It is holomorphic on a neighborhood
of $(u,v)=(0,0)$, including $\lambda=\mu$.

For $0\le k<r$, define the finite expression
\begin{equation}\label{urc:mixed-kernel}
 M_{r,k}(u;\lambda,\mu)
 =P_r(u)\lambda^{k-r-u}
  \sum_{j=0}^k\binom kj
      \frac{(\mu/\lambda-1)^j}{r-k+j+u}.
\end{equation}
For the critical degree use the removable germ
\begin{equation}\label{urc:critical-kernel}
 N_r(u;\epsilon,\lambda,\mu)
 =P_r(u)\left\{
 \epsilon^{-u}\lambda^{-u}
 \left(\frac1u+\sum_{j=1}^r\binom rj
                  \frac{(\mu/\lambda-1)^j}{u+j}\right)
 -\frac1u\right\}.
\end{equation}
At each fixed $\epsilon>0$ its apparent pole at $u=0$ cancels
before any differentiation. Its constant in $\log\epsilon$ is
$r!C_r(u;\lambda,\mu)$.

Finally let $\mathcal F_{r,t}(u,v)$ denote the holomorphic germ whose
coefficient dictionary is
\begin{equation}\label{urc:fp-germ}
 (-1)^{m+n}\partial_u^m\partial_v^n
          \mathcal F_{r,t}(0,0)=F_{m,n}^{r,t}.
\end{equation}
To specify it independently of a meromorphic regularization, set
\begin{align}
 R_{r,t}(u,v;x)={}&Z_r(u,x)Z_t(v,x)
 -P_r(u)P_t(v)x^{-r-t-2-u-v}\notag\\
 &-P_r(u)\sum_{k=0}^r a_{t,k}(v)x^{k-r-1-u}
 -P_t(v)\sum_{k=0}^t a_{r,k}(u)x^{k-t-1-v},
 \label{urc:remainder}
\end{align}
and define
\begin{align}
 \mathcal F_{r,t}(u,v)={}&\int_0^1R_{r,t}(u,v;x)\,dx
 -\frac{P_r(u)P_t(v)}{r+t+1+u+v}\notag\\
 &-P_r(u)\sum_{k=0}^{r-1}\frac{a_{t,k}(v)}{r-k+u}
 -P_t(v)\sum_{k=0}^{t-1}\frac{a_{r,k}(u)}{t-k+v}.
 \label{urc:fp-germ-construction}
\end{align}
Empty sums are zero. The critical terms $k=r$ and $k=t$ have no
added constant: every monomial $x^{-1}\log^h x$ has finite part
zero on $(0,1)$ in the coordinate \eqref{urc:finitepart}. This is
why retaining a whole meromorphic $1/u$ expression would give the
wrong Stieltjes coefficient convention.

\begin{theorem}[All singular terms in one joint germ]
\label{urc:full-theorem}
For fixed $r,t\ge0$ and $\lambda,\mu>0$, as $\epsilon\downarrow0$,
\begin{align}
 &\int_0^1Z_r(u,x+\lambda\epsilon)Z_t(v,x+\mu\epsilon)\,dx
 =\mathcal F_{r,t}(u,v)
   +\epsilon^{-r-t-1-u-v}B_{r,t}(u,v;\lambda,\mu)\notag\\
 &\quad+\sum_{k=0}^{r-1}a_{t,k}(v)
               \epsilon^{k-r-u}M_{r,k}(u;\lambda,\mu)
       +\sum_{k=0}^{t-1}a_{r,k}(u)
               \epsilon^{k-t-v}M_{t,k}(v;\mu,\lambda)\notag\\
 &\quad+a_{t,r}(v)N_r(u;\epsilon,\lambda,\mu)
       +a_{r,t}(u)N_t(v;\epsilon,\mu,\lambda)+o(1).
 \label{urc:full-expansion}
\end{align}
The remainder tends to zero locally uniformly for $(u,v)$ in a
complex neighborhood of zero. Every fixed mixed order derivative
may therefore be taken in this asymptotic formula. Applying
$(-1)^{m+n}\partial_u^m\partial_v^n$ and setting $u=v=0$ gives
the complete negative-power and logarithmic expansion of
\eqref{urc:integral}, through its constant.
\end{theorem}

The possible negative powers are especially sparse. They consist of
$\epsilon^{-(r+t+1)}$, together with $\epsilon^{-r},\ldots,
\epsilon^{-1}$ and $\epsilon^{-t},\ldots,\epsilon^{-1}$ when the
corresponding derivative orders are positive. The first power has
a logarithmic polynomial of degree at most $m+n$. The two finite
mixed families have degrees at most $m$ and $n$. The critical
terms give logarithmic polynomials of degrees at most $m+1$ and
$n+1$. Formula \eqref{urc:full-expansion} includes cancellations
between these polynomials automatically.

\begin{proof}
The Hurwitz recurrence gives \eqref{urc:local-decomposition} exactly.
Its analytic part has the Taylor expansion
\[
 A_t^{\mathrm{reg}}(v,z)=\sum_{k=0}^r a_{t,k}(v)z^k
                              +O(z^{r+1})
\]
uniformly for $v$ in a sufficiently small closed complex disk.
Expand the product of its singular and analytic parts. For the
singular--singular term, substitute $x=\epsilon y$ and complete
the interval to $(0,\infty)$. Its omitted tail tends locally
uniformly to $P_r(u)P_t(v)/(r+t+1+u+v)$, giving precisely the
beta term in \eqref{urc:full-expansion} and the first rational
constant in \eqref{urc:fp-germ-construction}.

For a mixed Taylor monomial the finite identity is
\begin{align}
 &P_r(u)\int_0^1
       (x+\lambda\epsilon)^{-r-1-u}(x+\mu\epsilon)^k\,dx
 \notag\\
 &\quad=P_r(u)\sum_{j=0}^k\binom kj
       ((\mu-\lambda)\epsilon)^j
 \frac{(\lambda\epsilon)^{k-j-r-u}
                  -(1+\lambda\epsilon)^{k-j-r-u}}
      {r-k+j+u}.
 \label{urc:exact-mixed-primitive}
\end{align}
It follows by expanding $x+\mu\epsilon=(x+\lambda\epsilon)
+(\mu-\lambda)\epsilon$ and integrating each power. At $k<r$,
the lower endpoint is $\epsilon^{k-r-u}M_{r,k}$, while the upper
endpoint tends to $-P_r(u)/(r-k+u)$. At $k=r$ the combination of
the lower endpoint and $-P_r(u)/u$ is exactly $N_r$.
The remaining upper-endpoint error is holomorphic at $u=0$ and
tends uniformly to zero; in particular,
$\{1-(1+\lambda\epsilon)^{-u}\}/u=O(\epsilon)$ uniformly near
zero. Exchange the two factors for the other mixed terms.

It remains to justify passage to the limit in the Taylor remainders.
For $|\Re u|,|\Re v|\le\eta<1/4$ and fixed positive rates, a
singular factor times the Taylor remainder is bounded by
$C(x+\epsilon)^{-\eta}$, hence by $Cx^{-\eta}$ on $(0,1)$.
This integrable majorant is uniform on compact parameter disks.
The product of analytic remainders is uniformly bounded. Dominated
convergence therefore gives the integral in
\eqref{urc:fp-germ-construction}, locally uniformly in $u,v$.
Cauchy's formula on slightly larger parameter disks gives the
same conclusion after every prescribed mixed order derivative.
This proves \eqref{urc:full-expansion}, including its joint rather
than merely raywise spectral meaning.

Coefficient extraction in \eqref{urc:fp-germ-construction} gives
\eqref{urc:fp-germ}: noncritical singular monomials have their
ordinary finite primitives at one, whereas every critical
$x^{-1}\log^h x$ contributes zero. Thus no exchange between an
unsubtracted divergent integral and a spectral derivative is used.
\end{proof}

\begin{proof}[Proof of Theorem~\ref{urc:constant-theorem}]
In \eqref{urc:full-expansion}, the beta term and the noncritical
mixed terms have strictly negative powers of $\epsilon$, also
after spectral differentiation. They cannot contribute to the
constant under discussion. In each critical term, retain the
constant in $\log\epsilon$, namely $r!C_r$ or $t!C_t$.
Equation \eqref{urc:analytic-coefficients} then gives
\eqref{urc:constant-formula}. At $u=0$,
\[
 C_r(0;\lambda,\mu)=(-1)^r\left\{-\log\lambda+
       \sum_{j=1}^r\binom rj\frac{(R-1)^j}{j}\right\}.
\]
The polynomial identity
\[
 \sum_{j=1}^r\binom rj\frac{(R-1)^j}{j}
       =\sum_{j=1}^r\frac{R^j-1}{j}
\]
follows by differentiating both sides in $R$ and using their common
value zero at $R=1$. This proves \eqref{urc:zeroth-contact}.
At $r=0$, expanding $(\lambda^{-u}-1)/u$ gives
\eqref{urc:no-argument-derivatives}.
\end{proof}

\subsection{Polylogarithmic coefficients of the divergent beta term}

The hypergeometric notation in \eqref{urc:beta-kernel} does not add
an unspecified family of constants. It is enough to differentiate
the base kernel in the shifts:
\[
 B_{r,t}(u,v;\lambda,\mu)
       =\partial_\lambda^r\partial_\mu^t
                      B_{0,0}(u,v;\lambda,\mu).
\]
For $\lambda>\mu$, put $d=\lambda-\mu$ and $z=d/\lambda$.
The substitution $y=d/(x+\lambda)$ gives
\begin{equation}\label{urc:polylog-beta}
 B_{0,0}(u,v;\lambda,\mu)
  =d^{-1-u-v}\int_0^z
       \frac{\exp\{(u+v)\log y-v\log(1-y)\}}{1-y}\,dy.
\end{equation}
Every coefficient of total spectral degree $N$ is a finite linear
combination of harmonic polylogarithms on the alphabet $\{0,1\}$
and logarithms, of weight at most $N+1$. This follows by shuffling
the logarithmic factors and integrating once with $dy/(1-y)$.
Shift differentiation preserves this finite function class. The
case $\mu>\lambda$ follows by exchange; the coincident case is
the elementary expression
\[
 B_{r,t}(u,v;\lambda,\lambda)
 =\frac{P_r(u)P_t(v)\lambda^{-r-t-1-u-v}}{r+t+1+u+v}.
\]
Thus every coefficient in the full subtraction belongs to an
explicit finite alphabet of harmonic numbers, logarithms,
polylogarithms, and Hurwitz or Stieltjes derivatives at one.

\subsection{Three completely explicit polygamma collision laws}

Write
\begin{equation}\label{urc:D}
 D(\lambda,\mu)=\frac{\log(\lambda/\mu)}{\lambda-\mu},
 \qquad D(\lambda,\lambda)=\frac1\lambda,
\end{equation}
and let $D_\lambda$ and $D_{\lambda\mu}$ be its indicated
derivatives. Explicitly, away from the diagonal,
\begin{align*}
 D_\lambda&=\frac{(\lambda-\mu)/\lambda-\log(\lambda/\mu)}
                  {(\lambda-\mu)^2},\\
 D_{\lambda\mu}
 &=\frac{(\lambda-\mu)(\lambda^{-1}+\mu^{-1})
                       -2\log(\lambda/\mu)}{(\lambda-\mu)^3}.
\end{align*}
Their diagonal values are $-1/(2\lambda^2)$ and
$1/(3\lambda^3)$. The inherited coincident evaluations \cite{LatticeIncoming} are
\[
 F_{0,0}^{0,0}=2\gamma_1-2\zeta(2),\qquad
 F_{0,0}^{1,0}=\zeta(2),\qquad
 F_{0,0}^{1,1}=2\zeta(3)+4\zeta'(3).
\]
Using these established values in the new collision formula gives
\begin{align}
 \int_0^1\psi(x+\lambda\epsilon)\psi(x+\mu\epsilon)\,dx
 ={}&\frac{D(\lambda,\mu)}\epsilon-2\gamma\log\epsilon
       +2\gamma_1-2\zeta(2)-\gamma\log(\lambda\mu)+o(1),
 \label{urc:digamma-example}\\
 \int_0^1\psi'(x+\lambda\epsilon)\psi(x+\mu\epsilon)\,dx
 ={}&\frac{D_\lambda(\lambda,\mu)}{\epsilon^2}
       -\frac{\gamma}{\lambda\epsilon}
       +\boxed{\zeta(2)\left\{\frac\mu\lambda+
                               \log\frac\mu\lambda\right\}}
       +o(1),
 \label{urc:ordered-example}\\
 \int_0^1\psi'(x+\lambda\epsilon)\psi'(x+\mu\epsilon)\,dx
 ={}&\frac{D_{\lambda\mu}(\lambda,\mu)}{\epsilon^3}
       +\frac{\zeta(2)}\epsilon
                       \left(\frac1\lambda+\frac1\mu\right)
       +4\zeta(3)\log\epsilon\notag\\
 &+4\zeta'(3)+2\zeta(3)
       \left\{3+\log(\lambda\mu)-\frac\mu\lambda-
                                      \frac\lambda\mu\right\}+o(1).
 \label{urc:trigamma-example}
\end{align}
The ordered law \eqref{urc:ordered-example} is invariant under a
common scaling of the two rates. Its exchange companion has
constant $\zeta(2)\{\lambda/\mu+\log(\lambda/\mu)\}$. Their
sum equals $\zeta(2)(\lambda/\mu+\mu/\lambda)$, in exact
agreement with the endpoint derivative of the digamma product.
This independently checks the logarithmic signs and the ratio
terms; exchange symmetry alone would not determine either ordered
constant separately.

\subsection{A first Stieltjes collision with a dilogarithm}

Let $\gamma_j=\gamma_j(1)$ and put
\begin{equation}\label{urc:E}
 E(\lambda,\mu)=
 \frac{\log(\lambda/\mu)\log\lambda+
               \operatorname{Li}_2(1-\mu/\lambda)}{\lambda-\mu},
 \qquad E(\lambda,\lambda)=\frac{1+\log\lambda}{\lambda}.
\end{equation}
This is the ordinary convergent integral of
$\log(y+\lambda)/\{(y+\lambda)(y+\mu)\}$ on $(0,\infty)$,
as follows either from \eqref{urc:polylog-beta} or one integration
by parts. The standard coincident product germ gives
\begin{equation}\label{urc:F10}
 F_{1,0}^{0,0}=\frac32\gamma_2+2\gamma\zeta(2)-\zeta(3).
\end{equation}
Consequently the complete first-index expansion is
\begin{align}
 \int_0^1\gamma_1(x+\lambda\epsilon)
                    \gamma_0(x+\mu\epsilon)\,dx
 ={}&\frac{D(\lambda,\mu)\log\epsilon+E(\lambda,\mu)}\epsilon
       -\frac\gamma2(\log\epsilon)^2\notag\\
 &-(\gamma\log\lambda+\gamma_1)\log\epsilon
       +\frac32\gamma_2+2\gamma\zeta(2)-\zeta(3)\notag\\
 &-\frac\gamma2(\log\lambda)^2-\gamma_1\log\mu+o(1).
 \label{urc:first-Stieltjes-example}
\end{align}
All logarithms and the dilogarithm in this formula are real for
positive rates, with the displayed continuous diagonal limit.

For completeness, \eqref{urc:F10} can be extracted without a
numerical relation search. Put $w=u+v$,
$Q=\Gamma(1+w)/\{\Gamma(1+u)\Gamma(1+v)\}$ and
$E_0(u)=\pi u/\sin\pi u$. The inherited coincident Hurwitz
product formula gives the holomorphic finite-part germ
\begin{align*}
 \mathcal F_{0,0}(u,v)
 ={}&\frac{w-Q\{uE_0(v)+vE_0(u)\}}{uvw}\\
 &+\frac{Z(v,1)-QE_0(u)Z(w,1)}u
  +\frac{Z(u,1)-QE_0(v)Z(w,1)}v.
\end{align*}
Each divided difference is removable. Its degree-zero term is
$2\gamma_1-2\zeta(2)$ and its coefficient of $u$ is
$\zeta(3)-3\gamma_2/2-2\gamma\zeta(2)$. Applying the sign
convention \eqref{urc:stieltjes} proves \eqref{urc:F10}.

\subsection{Verification, scope, and further questions}

The accompanying verifier checks the finite mixed primitive,
the rising-factorial cancellations, the compact ratio polynomial,
and contact jets by two exact coefficient constructions. Separate
high-precision ordinary quadratures check the three beta kernels
and the three polygamma collisions at diagonal, unequal integral,
and unequal rational rates. The numerical collision residuals are
compared at decreasing positive $\epsilon$ after the complete
displayed subtraction. They test the asymptotic constants and
normalization; they are not certified error enclosures.

The result concerns fixed positive rates. Paths with a vanishing or
unbounded ratio require an additional asymptotic regime; inserting
a rate that depends on $\epsilon$ into a fixed-compactness proof
would not justify that regime. The following are concrete next
questions.
\begin{enumerate}
\item Determine the complete transition for
$a=\lambda\epsilon^\alpha$, $b=\mu\epsilon^\beta$ with
$\alpha\ne\beta$, including the boundary between the two
scales. The explicit finite ratio polynomial and the beta kernel
make every potential extra power visible.
\item Extend the locality theorem to three independently shifted
Stieltjes factors. With more than two factors, products of two
singular parts and one analytic Taylor coefficient can already
affect the constant; the required kernel should be stated before
any reduction to a smaller alphabet is claimed.
\item Classify linear combinations of ordered derivative products
whose collision constant depends only on the ratio $\mu/\lambda$,
or is independent of the rates. Formula \eqref{urc:finite-coefficients}
reduces this to an exact finite polynomial problem at each index.
\item Continue the rates to a common complex sector and derive
the jumps when the lower-endpoint logarithms cross their chosen
cuts. The local beta integral supplies a natural branch convention.
\end{enumerate}
